\documentclass[10pt,oneside]{book}

\usepackage[T1]{fontenc}
\usepackage[utf8]{inputenc}
\usepackage[french]{babel}
\usepackage{amsmath,amssymb,mathrsfs}
\usepackage{longtable}
\usepackage[all]{xy}
\usepackage{tikz-cd}
\usepackage[a4paper,margin=22mm]{geometry}
\usepackage[
  colorlinks=true,
  linkcolor=blue,
  hypertexnames=false,
  unicode=true,
  pdfencoding=auto
]{hyperref}
\input{user/Documents/ega_scan_check_20261101/ega-fr/source/preamble-base.tex}
\hypersetup{
  pdftitle={Éléments de géométrie algébrique -- édition française cumulative},
  pdfauthor={Alexander Grothendieck et Jean Dieudonné},
  pdfsubject={Texte diplomatique français cumulatif depuis l'autorité NUMDAM}
}

\renewcommand{\thesection}{\arabic{section}}
\renewcommand{\thesubsection}{\thesection.\arabic{subsection}}
\newcommand{\oldpage}[2][]{%
  \setcounter{footnote}{0}%
  \marginpar{\scriptsize\textbf{#1}~|~#2}\ignorespaces}
\newenvironment{env}[1][]{\par\medskip\noindent\textbf{(#1)}\ }{\par}
\newenvironment{definition}[1][]{%
  \par\medskip\noindent\itshape Définition (#1). ---\ }{\par}
\newenvironment{proposition}[1][]{%
  \par\medskip\noindent\itshape Proposition (#1). ---\ }{\par}
\newenvironment{theorem}[1][]{%
  \par\medskip\noindent\itshape Théorème (#1). ---\ }{\par}
\newenvironment{corollary}[1][]{%
  \par\medskip\noindent\itshape Corollaire (#1). ---\ }{\par}
\newenvironment{remark}[1][]{%
  \par\medskip\noindent\itshape Remarque (#1). ---\ }{\par}
\newenvironment{remarks}[1][]{%
  \par\medskip\noindent\itshape Remarques (#1). ---\ }{\par}
\newenvironment{scholium}[1][]{%
  \par\medskip\noindent\textit{Scholie (#1). ---\ }}{\par}
\newenvironment{lemma}[1][]{%
  \par\medskip\noindent\itshape Lemme (#1). ---\ }{\par}
\newenvironment{example}[1][]{%
  \par\medskip\noindent\textit{Exemple (#1). ---\ }}{\par}
\newenvironment{examples}[1][]{%
  \par\medskip\noindent\textit{Exemples (#1). ---\ }}{\par}
\newenvironment{notation}[1][]{%
  \par\medskip\noindent\textit{Notations (#1). ---\ }}{\par}
\newenvironment{proof}{%
  \par\medskip\noindent\textit{Démonstration. ---\ }}{\par}
\newenvironment{namedtheorem}[2][]{%
  \par\medskip\noindent\itshape #2 (#1). ---\ }{\par}
\newenvironment{namedlemma}[2][]{%
  \par\medskip\noindent\itshape #2 (#1). ---\ }{\par}

\begin{document}
\setcounter{section}{4}\setcounter{subsection}{6}
\begin{proposition}[4.6.7]
\label{III.4.6.7-fr}
Soient $Z$ un préschéma localement noethérien, $X$, $Y$ deux $Z$-préschémas
tels que les morphismes structuraux $g:X\to Z$, $h:Y\to Z$ soient propres.
Soient $f:X\to Y$ un $Z$-morphisme, $z$ un point de $Z$, et soit
\[
f_z=f\times_Z1:
X\otimes_Zk(z)\longrightarrow Y\otimes_Zk(z).
\]
\begin{enumerate}
\item[(i)] Si $f_z$ est un morphisme fini (resp. une immersion fermée), il
existe un voisinage ouvert $U$ de $z$ tel que le morphisme
$g^{-1}(U)\to h^{-1}(U)$, restriction de $f$, soit un morphisme fini (resp.
une immersion fermée).
\item[(ii)] On suppose de plus que $g$ soit un morphisme plat. Alors, si
$f_z$ est un isomorphisme, il existe un voisinage ouvert $U$ de $z$ tel que
le morphisme $g^{-1}(U)\to h^{-1}(U)$, restriction de $f$, soit un
isomorphisme.
\end{enumerate}
\end{proposition}

\begin{proof}
Dans les deux cas, il suffira de prouver que pour tout $y\in h^{-1}(z)$ il
existe un voisinage $V_y$ de $y$ tel que la restriction
$f^{-1}(V_y)\to V_y$ de $f$ soit un morphisme fini (resp. une immersion
fermée, un isomorphisme)~; il en résultera alors en effet que si $V$ est la
réunion des $V_y$, la restriction $f^{-1}(V)\to V$ de $f$ est un morphisme
fini (resp. une immersion fermée, un isomorphisme) (II, 6.1.1 et I, 4.2.4).
Comme $h$ est un morphisme fermé, il existera un voisinage ouvert $U$ de $z$
tel que $h^{-1}(U)\subset V$, et la proposition sera démontrée.

\noindent\textit{(i)} Notons tout d'abord que $f$ est un morphisme propre
(II, 5.4.3)~; si on
\oldpage[III]{143}
suppose $f_z$ fini, l'existence pour tout
$y\in h^{-1}(z)$ d'un voisinage $V_y$ tel que $f^{-1}(V_y)\to V_y$ soit fini
résulte de (\hyperref[III.4.4.11-fr]{4.4.11}).
Pour traiter le cas où $f_z$ est une immersion fermée, on peut donc déjà
supposer que le morphisme $f$ est \emph{fini}, donc que
$X=\operatorname{Spec}(\mathcal B)$, où $\mathcal B$ est une
$\mathcal O_Y$-Algèbre cohérente, le morphisme $f$ correspondant (\textbf{II}, 1.2.7)
à l'homomorphisme canonique $u:\mathcal O_Y\to\mathcal B$. Si l'on prouve
que pour tout $y\in h^{-1}(z)$, l'homomorphisme
$u_y:\mathcal O_y\to\mathcal B_y$ est surjectif, il en résultera que pour un
voisinage $V_y$ de $y$, $u|V_y$ sera surjectif, le faisceau
$\operatorname{Coker}u$ étant cohérent (\textbf{0}\textsubscript{I}, 5.3.4 et 5.2.2).
Cela étant, le morphisme fini $f_z$ correspond à l'homomorphisme
$v=u\otimes1:
\mathcal O_Y\otimes_{\mathcal O_Z}\boldsymbol{k}(z)
\longrightarrow
\mathcal B\otimes_{\mathcal O_Z}\boldsymbol{k}(z)$,
et l'hypothèse que $f_z$ est une immersion fermée entraîne que
l'homomorphisme
$u_y\otimes1:
\mathcal O_y\otimes_{\mathcal O_z}\boldsymbol{k}(z)
\longrightarrow
\mathcal B_y\otimes_{\mathcal O_z}\boldsymbol{k}(z)$
est surjectif. Comme $\mathcal B_y$ est un $\mathcal O_y$-module de type
fini et $\mathcal O_y$ un anneau local noethérien, la conclusion résulte
comme dans (\hyperref[III.4.6.3-fr]{4.6.3}) du lemme de Nakayama.

\noindent (ii) Le même raisonnement que ci-dessus montre qu'il
suffit cette fois de prouver que $u_y$ est bijectif, sachant que
$u_y\otimes1$ est bijectif.

Cela résultera du lemme suivant~:
\end{proof}

\begin{lemma}[4.6.7.1]
\label{III.4.6.7.1-fr}
Soient $A$, $B$ deux anneaux locaux noethériens,
$\rho:A\to B$ un homomorphisme local, $u:N\to M$ un homomorphisme de
$B$-modules. On suppose que $M$ est un $A$-module plat, $N$ un $B$-module
de type fini et que
$u\otimes1:N\otimes_Ak\longrightarrow M\otimes_Ak$
(où $k$ est le corps résiduel de $A$) est injectif. Alors $N$ est un
$A$-module plat et $u$ est injectif.
\end{lemma}

\begin{proof}
Pour établir la première assertion, il faut montrer que pour tout couple de
$A$-modules \emph{de type fini} $P$, $Q$ et tout $A$-homomorphisme injectif
$v:P\to Q$,
$1_N\otimes v:N\otimes_AP\to N\otimes_AQ$ est injectif. Or, on a le
diagramme commutatif
\[
\xymatrix{
N\otimes_AP\ar[r]^{1_N\otimes v}\ar[d]_{u\otimes1_P}
  &N\otimes_AQ\ar[d]^{u\otimes1_Q}\\
M\otimes_AP\ar[r]_{1_M\otimes v}
  &M\otimes_AQ
}
\]
et comme $1_M\otimes v$ est injectif par hypothèse, il suffira de prouver
qu'il en est de même de $u\otimes1_P$. Soit $\mathfrak m$ l'idéal maximal
de $A$~; la filtration $\mathfrak m$-adique sur le $A$-module
$N\otimes_AP$ est aussi sa filtration $\mathfrak mB$-adique en tant que
$B$-module~; la topologie définie par cette filtration est donc \emph{séparée},
puisque $B$ est noethérien, que $\mathfrak mB$ est contenu dans le radical de
$B$, et que $N\otimes_AP$ est un $B$-module de type fini, $N$ étant un
$B$-module de type fini et $P$ un $A$-module de type fini
(\textbf{0}\textsubscript{I}, 7.3.5). Il suffit donc de prouver que l'homomorphisme
$\operatorname{gr}(u\otimes1_P):
\operatorname{gr}_\bullet(N\otimes_AP)
\longrightarrow
\operatorname{gr}_\bullet(M\otimes_AP)$
(où les modules gradués sont relatifs aux filtrations $\mathfrak m$-adiques)
est injectif (Bourbaki, \emph{Alg. comm.}, chap. III, \S\,2, n\textsuperscript{o}~8, cor. 1 du
th. 1). Notons maintenant que puisque $M$ est un $A$-module plat, les
homomorphismes
$M\otimes_A(\mathfrak m^nP)
\longrightarrow
\mathfrak m^n(M\otimes_AP)$
sont bijectifs~; il en est donc de même de l'homomorphisme canonique
\[
\varphi_M:
\operatorname{gr}_0(M)\otimes_A\operatorname{gr}_\bullet(P)
\longrightarrow
\operatorname{gr}_\bullet(M\otimes_AP).
\]

\oldpage[III]{144}

Or, on a un diagramme commutatif
\[
\xymatrix{
\operatorname{gr}_0(N)\otimes_A\operatorname{gr}_\bullet(P)
  \ar[r]^{\operatorname{gr}_0(u)\otimes1}\ar[d]_{\varphi_N}
  &\operatorname{gr}_0(M)\otimes_A\operatorname{gr}_\bullet(P)
  \ar[d]^{\varphi_M}\\
\operatorname{gr}_\bullet(N\otimes_AP)
  \ar[r]_{\operatorname{gr}(u\otimes1)}
  &\operatorname{gr}_\bullet(M\otimes_AP)
}
\]
dans lequel $\varphi_M$ est bijectif, $\varphi_N$ surjectif~; en outre,
$\operatorname{gr}_0(u)$ est injectif par hypothèse, et comme
$\operatorname{gr}_0(N)\otimes_A\operatorname{gr}_\bullet(P)
=\operatorname{gr}_0(N)\otimes_k\operatorname{gr}_\bullet(P)$,
$\operatorname{gr}_0(M)\otimes_A\operatorname{gr}_\bullet(P)
=\operatorname{gr}_0(M)\otimes_k\operatorname{gr}_\bullet(P)$,
$\operatorname{gr}_0(u)\otimes1$ est aussi injectif. On en conclut que
$\operatorname{gr}(u\otimes1)$ est injectif, ce qui achève de démontrer la
première assertion. La seconde se déduit du raisonnement précédent en faisant
$P=A$.
\end{proof}

\begin{proposition}[4.6.8]
\label{III.4.6.8-fr}
Soient $Z$ un préschéma localement noethérien, $X$, $Y$, deux $Z$-préschémas
tels que les morphismes structuraux $g:X\to Z$, $h:Y\to Z$ soient propres,
$Z'$ une partie fermée de $Z$, $X'=g^{-1}(Z')$, $Y'=h^{-1}(Z')$ ses images
réciproques,
$\widehat X=X_{/X'}$, $\widehat Y=Y_{/Y'}$, $\widehat Z=Z_{/Z'}$
les complétés formels de $X$, $Y$, $Z$ le long de ces parties fermées,
$f:X\to Y$ un $Z$-morphisme, $\widehat f:\widehat X\to\widehat Y$ son
prolongement aux complétés. Pour que $\widehat f$ soit un isomorphisme (\textup{resp.}
une immersion fermée), il faut et il suffit qu'il existe un voisinage ouvert
$U$ de $Z'$ tel que le morphisme $g^{-1}(U)\to h^{-1}(U)$, restriction de
$f$, soit un isomorphisme (\textup{resp.} une immersion fermée).
\end{proposition}

\begin{proof}
La suffisance de la condition est immédiate (\textbf{I}, 10.14.7). Pour en montrer la
nécessité, il suffit encore de prouver que pour tout $y\in Y'$, il existe un
voisinage ouvert $V_y$ de $y$ tel que la restriction
$f^{-1}(V_y)\to V_y$ de $f$ soit un isomorphisme (resp. une immersion
fermée), par le même raisonnement que dans
(\hyperref[III.4.6.7-fr]{4.6.7}). On est ainsi ramené au cas où $Y=Z$,
$Y=\operatorname{Spec}(A)$ étant affine noethérien. Par hypothèse
(I, 10.9.1 et 10.14.2) la fibre $f^{-1}(y)$ est réduite à un point pour
$y\in Y'$, donc comme $f$ est propre (\textbf{II}, 5.4.3), il existe un voisinage
ouvert $U$ de $y$ tel que la restriction $f^{-1}(U)\to U$ de $f$ soit un
morphisme \emph{fini} (\hyperref[III.4.4.11-fr]{4.4.11}). On peut donc déjà supposer
que $f$ soit un morphisme fini, donc $X=\operatorname{Spec}(B)$, où $B$ est
une $A$-algèbre \emph{finie} sur $A$. Si $Y'=V(\mathfrak J)$, on a alors
$\widehat Y=\operatorname{Spf}(\widehat A)$, $\widehat X=\operatorname{Spf}(\widehat B)$,
$\widehat A$ étant le séparé complété de $A$ pour la topologie
$\mathfrak J$-préadique, $\widehat B$ le séparé complété de $B$ pour la
topologie $\mathfrak JB$-préadique, ou (ce qui revient au même), le séparé
complété du $A$-module $B$ pour la topologie $\mathfrak J$-préadique~; en
outre, $\widehat f$ est le morphisme de schémas formels affines correspondant
au prolongement continu $\widehat\varphi:\widehat A\to\widehat B$ de
l'homomorphisme canonique d'anneaux $\varphi:A\to B$, et l'hypothèse est que
$\widehat\varphi$ est surjectif (resp. bijectif) (\textbf{I}, 10.14.2). Or,
$\widehat\varphi$ est aussi le prolongement continu de $\varphi$ considéré
comme homomorphisme de $A$-modules~; on sait alors (\textbf{I}, 10.8.14) qu'il existe
un voisinage ouvert $U$ de $Y'$ tel que la restriction à $U$ de
l'homomorphisme $\widetilde\varphi:\widetilde A\to\widetilde B$ de
$\mathcal O_Y$-Modules soit surjectif (resp. bijectif), ce qui achève la
démonstration.
\end{proof}

\oldpage[III]{145}

\subsection{Un critère d'amplitude}

\begin{theorem}[4.7.1]
\label{III.4.7.1-fr}
Soient $Y$ un préschéma localement noethérien, $f:X\to Y$ un morphisme
propre, $\mathcal L$ un $\mathcal O_X$-Module inversible, $y$ un point de
$Y$, $X_y=X\otimes_Y\boldsymbol{k}(y)=f^{-1}(y)$, $g$ la projection de $X_y$ dans $X$.
Si $\mathcal L_y=g^*(\mathcal L)=\mathcal L\otimes_{\mathcal O_Y}\boldsymbol{k}(y)$ est
ample sur $X_y$, il existe un voisinage ouvert $U$ de $y$ dans $Y$ tel que
$\mathcal L|f^{-1}(U)$ soit ample pour la restriction de $f$ à $f^{-1}(U)$.
\end{theorem}

\begin{proof}
\noindent I) Posons
$Y'=\operatorname{Spec}(\mathcal O_y)$, $X'=X\times_YY'$, et soit
$\mathcal L'=\mathcal L\otimes_{\mathcal O_Y}\mathcal O_y$~;
nous allons d'abord prouver que $\mathcal L'$ est \emph{ample pour}
$f'=f_{(Y')}$. On a le diagramme commutatif
\[
\xymatrix{
X\ar[d]_f & X'\ar[l]\ar[d]_{f'} & X_y\ar[l]\ar[d]_{f_y}\\
Y & Y'\ar[l] & \operatorname{Spec}(\boldsymbol{k}(y))\ar[l]
}
\]
Comme $f'$ est propre (\textbf{II}, 5.4.2, (iii)) et $\mathcal O_y$ noethérien, on
voit qu'on peut se borner au cas où
$Y=Y'=\operatorname{Spec}(\mathcal O_y)$, donc $X=X'$, supposer que
$\mathcal L_y$ est ample pour $f_y$ et prouver que $\mathcal L$ est ample
pour $f$ (\textbf{II}, 4.6.6). Nous allons appliquer le critère
(\hyperref[III.2.6.1-fr]{2.6.1, \emph{c)}}) et montrer en fait que pour tout
$\mathcal O_X$-Module cohérent $\mathcal F$, il existe un entier $N$ tel que
$H^1(X,\mathcal F(n))=0$ pour tout $n\geq N$, avec
$\mathcal F(n)=\mathcal F\otimes\mathcal L^{\otimes n}$.
Notons que $y$ est un point fermé de $Y$ correspondant à l'idéal maximal
$\mathfrak m$ de $\mathcal O_y$~; $X_y$ est donc un sous-préschéma fermé de
$X$ défini \textbf{par} l'Idéal cohérent
$\mathcal J=f^*(\widetilde{\mathfrak m})\mathcal O_X=
\mathfrak m\mathcal O_X$ de $\mathcal O_X$ (\textbf{I}, 4.4.5), et $g$ l'injection
canonique. Considérons alors la $\boldsymbol{k}(y)$-algèbre graduée
$S=\operatorname{gr}(\mathcal O_y)
=\bigoplus_{j\geq0}\mathfrak m^j/\mathfrak m^{j+1}$,
qui est de type fini puisque $\mathcal O_y$ est noethérien~; la
$\mathcal O_X$-Algèbre $\mathcal S=f^*(\widetilde S)$ est donc
quasi-cohérente et de type fini, et elle est évidemment annulée par
$\mathcal J$, donc si on pose $\mathcal S_y=g^*(\mathcal S)$,
$\mathcal S_y$ est une $\mathcal O_{X_y}$-Algèbre quasi-cohérente de type
fini, et $\mathcal S=g_*(\mathcal S_y)$. Posons d'autre part,
$\mathcal M_j=\mathfrak m^j\mathcal F/\mathfrak m^{j+1}\mathcal F$ et
$\mathcal M=\bigoplus_{j\geq0}\mathcal M_j
=\operatorname{gr}(\mathcal F)$~;
comme $\mathcal F$ est cohérent, $\mathcal M$ est un $\mathcal S$-Module
quasi-cohérent de type fini (\textbf{0}, 10.1.1) qui est aussi annulé par
$\mathcal J$, de sorte que si l'on pose
$\mathcal M'_j=g^*(\mathcal M_j)$, $\mathcal M'=g^*(\mathcal M)=\bigoplus_{j\geq0}\mathcal M'_j$
est un $\mathcal S_y$-Module gradué quasi-cohérent de type
fini tel que $\mathcal M=g_*(\mathcal M')$. En outre, si on pose
$\mathcal M'_j(n)=\mathcal M'_j\otimes\mathcal L_y^{\otimes n}$, on a
$\mathcal M'_j(n)=g^*((\mathcal M_j(n)))$. Cela étant, $f_y$ est propre
(\textbf{II}, 5.4.2, (iii)) et $\mathcal L_y$ est ample, donc $f_y$ est \emph{projectif}
(\textbf{II}, 5.5.4 et 4.6.11), et on peut appliquer à
$\operatorname{Spec}(\boldsymbol{k}(y))$, $f_y$, $\mathcal S_y$, $\mathcal L_y$ et
$\mathcal M'$ le théorème (\hyperref[III.2.4.1-fr]{2.4.1, (ii)})~: il
existe un entier $N$ tel que pour $n\geq N$, on ait
$H^q(X_y,\mathcal M'_j(n))=0$ pour tout $q>0$ \emph{et tout} $j$~; par suite, on a
aussi $H^q(X,\mathcal M_j(n))=0$ pour tout $q>0$ et tout $j$
(G, II, 4.9.1). Posons alors
$\mathcal F(n)_j=\mathcal F(n)/\mathfrak m^{j+1}\mathcal F(n)$,
de sorte que $\mathcal F(n)_{j-1}=\mathcal F(n)_j/\mathcal M_j(n)$ pour
$j\geq1$ et $\mathcal F(n)_0=\mathcal M_0(n)$. On a
$H^1(X,\mathcal F(n)_0)=0$, et, par la suite exacte de cohomologie,
$H^1(X,\mathcal F(n)_j)=H^1(X,\mathcal F(n)_{j-1})$ pour tout $j\geq1$,
donc $H^1(X,\mathcal F(n)_j)=0$ pour tout $j\geq0$. On conclut donc de
(\hyperref[III.4.2.1-fr]{4.2.1}) que $H^1(X,\mathcal F(n))=0$, ce qui
achève de prouver notre assertion.

\noindent II) Revenons aux notations du début de la démonstration,
et remarquons qu'on
\oldpage[III]{146}
peut toujours supposer $Y=\operatorname{Spec}(A)$
affine~; comme $f'$ est de type fini et $\mathcal L'$ ample pour $f'$, il
existe un entier $m>0$ tel que $\mathcal L'^{\otimes m}$ soit très ample pour
$f'$ (\textbf{I}I, 4.6.11)~; remplaçant au besoin $\mathcal L$ par
$\mathcal L^{\otimes m}$, on peut se borner à considérer le cas où
$\mathcal L'$ est \emph{très ample} pour $f'$, et à prouver que
$\mathcal L|f^{-1}(U)$ est alors très ample pour $f$. Comme $f'$ est propre,
il existe alors une $Y'$-immersion \emph{fermée}
$j:X'\to P=\mathbf P_{Y'}^r$ pour un entier $r>0$ convenable, telle que
$\mathcal L'$ soit isomorphe à $j^*(\mathcal O_P(1))$ (\textbf{II}, 5.5.4, (ii))~;
cette immersion correspond canoniquement à un $\mathcal O_{X'}$-homomorphisme
\emph{surjectif}
$u:\mathcal O_{X'}^{r+1}\longrightarrow\mathcal L'$
(\textbf{II}, 4.2.3). Ce dernier correspond (\textbf{0}\textsubscript{I}, 5.1.1) à la donnée
de $r+1$ sections $s'_i$ ($0\leq i\leq r$) de $\mathcal L'$ au-dessus de
$X'$ qui engendrent ce $\mathcal O_{X'}$-Module. Ces sections sont aussi par
définition des sections de $f'_*(\mathcal L')$ au-dessus de $Y'$~; on a
$f'_*(\mathcal L')=f_*(\mathcal L)\otimes_{\mathcal O_Y}\mathcal O_{Y'}$
(\textbf{0}\textsubscript{I}, 5.4.10), $Y$ est affine et $\mathcal O_y$ est l'anneau
local en l'idéal premier $\mathfrak j_y$ de $A$, donc on a
$s'_i=s''_i/t_i$, où les $s''_i$ sont des sections de $f_*(\mathcal L)$
au-dessus de $Y$ et les $t_i$ des éléments de $A$ n'appartenant pas à
$\mathfrak j_y$~; on en conclut qu'il existe un voisinage ouvert affine $V$
de $y$ dans $Y$ et des sections $s_i$ de $f_*(\mathcal L)|V$ telles que
$s'_i=s_i/1$ (on rappelle que l'espace $Y'$ est contenu dans $V$, cf. \textbf{I},
2.4.2). Les $s_i$ sont alors des sections de $\mathcal L$ au-dessus de
$f^{-1}(V)$, définissant donc un homomorphisme
$v:\bigl(\mathcal O_X|f^{-1}(V)\bigr)^{r+1}
\longrightarrow\mathcal L|f^{-1}(V)$
qui, par hypothèse, est surjectif en tous les points de $f^{-1}(y)$~; comme
$\operatorname{Coker}(v)$ est cohérent (\textbf{0}\textsubscript{I}, 5.3.4), son
support est fermé (\textbf{0}\textsubscript{I}, 5.2.2) et par suite il existe un
voisinage ouvert $W\subset f^{-1}(V)$ de $f^{-1}(y)$ tel que la restriction
de $v$ à $W$ soit un homomorphisme \emph{surjectif}. Puisque le morphisme $f$ est
fermé, on peut supposer que $W$ est de la forme $f^{-1}(U)$, où $U$ est un
voisinage ouvert de $y$, et la conclusion résulte alors de (\textbf{II}, 4.2.3).
\end{proof}

\subsection{Morphismes finis de préschémas formels}

\begin{proposition}[4.8.1]
\label{III.4.8.1-fr}
Soient $\mathfrak Y$ un préschéma formel localement noethérien, $\mathcal K$ un
Idéal de définition de $\mathfrak Y$, $f:\mathfrak X\to\mathfrak Y$ un
morphisme de préschémas formels. Les conditions suivantes sont équivalentes~:
\begin{enumerate}
\item[a)] $\mathfrak X$ est localement noethérien, $f$ est un morphisme adique
  \textup{(\textbf{I}, 10.12.1)} et si l'on pose
  $\mathcal J=f^*(\mathcal K)\mathcal O_{\mathfrak X}$, le morphisme
  $f_0:(\mathfrak X,\mathcal O_{\mathfrak X}/\mathcal J)
  \longrightarrow
  (\mathfrak Y,\mathcal O_{\mathfrak Y}/\mathcal K)$
  déduit de $f$ est fini.
\item[b)] $\mathfrak X$ est localement noethérien et est limite inductive d'un
  $(Y_n)$-système inductif adique $(X_n)$ tel que le morphisme
  $X_0\to Y_0$ soit fini.
\item[c)] Tout point de $\mathfrak Y$ possède un voisinage ouvert formel affine
  noethérien $V$ tel que $f^{-1}(V)$ soit un ouvert formel affine et que
  $\Gamma(f^{-1}(V),\mathcal O_{\mathfrak X})$ soit un
  $\Gamma(V,\mathcal O_{\mathfrak Y})$-module de type fini.
\end{enumerate}
\end{proposition}

\begin{proof}
Il est immédiat que \emph{a)} entraîne \emph{b)} en vertu de (\textbf{I}, 10.12.3). Pour voir que \emph{b)}
entraîne \emph{c)}, on peut supposer que
$\mathfrak Y=\operatorname{Spf}(B)$, où $B$ est adique noethérien et
$\mathcal K=\mathfrak R^\Delta$, où $\mathfrak R$ est un idéal de définition
de $B$. Par hypothèse, $X_0$ est un schéma affine dont l'anneau $A_0$ est un
$B/\mathfrak R$-module de type fini (\textbf{II}, 6.1.3). En vertu de (\textbf{I}, 5.1.9), chacun
des $X_n$ est un schéma affine, et si $A_n$ est son anneau, l'hypothèse \emph{b)}
entraîne que pour $m\leq n$, $A_m$ est isomorphe à
$A_n/\mathfrak R^{m+1}A_n$. On en déduit que $\mathfrak X$ est isomorphe à
$\operatorname{Spf}(A)$, où $A=\varprojlim_n A_n$~; on conclut en vertu de
(\textbf{0}\textsubscript{I}, 7.2.9). Enfin, pour prouver que \emph{c)}
\oldpage[III]{147}
entraîne \emph{a)}, on peut se borner encore au cas
$\mathfrak Y=\operatorname{Spf}(B)$, $\mathfrak X=\operatorname{Spf}(A)$,
$A$ étant une $B$-algèbre finie~; comme $A/\mathfrak RA$ est alors une
$B/\mathfrak R$-algèbre finie, il résulte de (\textbf{I}, 10.10.9) que les conditions
de \emph{a)} sont satisfaites.
\end{proof}

\begin{definition}[4.8.2]
\label{III.4.8.2-fr}
Lorsque les propriétés équivalentes \textup{a), b), c)} de
\textup{(\hyperref[III.4.8.1-fr]{4.8.1})} sont vérifiées, on dit que le morphisme $f$
est fini, ou que $\mathfrak X$ est un $\mathfrak Y$-préschéma
formel fini, ou un préschéma formel fini au-dessus de $\mathfrak Y$.
\end{definition}

\begin{proposition}[4.8.3]
\label{III.4.8.3-fr}
\begin{enumerate}
\item[\textup{(i)}] Une immersion fermée de préschémas formels localement noethériens
  est un morphisme fini.
\item[\textup{(ii)}] Le composé de deux morphismes finis de préschémas formels
  localement noethériens est un morphisme fini.
\item[\textup{(iii)}] Soient $\mathfrak X$, $\mathfrak Y$, $\mathfrak S$ trois
  préschémas formels localement noethériens, $f:\mathfrak X\to\mathfrak S$ un
  morphisme fini, $g:\mathfrak Y\to\mathfrak S$ un morphisme~; alors le
  morphisme
  $\mathfrak X\times_{\mathfrak S}\mathfrak Y\to\mathfrak Y$ est fini.
\item[\textup{(iv)}] Soient $\mathfrak S$ un préschéma formel localement noethérien,
  $\mathfrak X'$, $\mathfrak Y'$ deux préschémas formels localement
  noethériens tels que $\mathfrak X'\times_{\mathfrak S}\mathfrak Y'$ soit
  localement noethérien. Si $\mathfrak X$, $\mathfrak Y$ sont des
  $\mathfrak S$-préschémas formels localement noethériens,
  $f:\mathfrak X\to\mathfrak X'$, $g:\mathfrak Y\to\mathfrak Y'$ deux
  $S$-morphismes finis, alors $f\times_{\mathfrak S}g$ est un
  morphisme fini.
\item[\textup{(v)}] Soient $f:\mathfrak X\to\mathfrak Y$,
  $g:\mathfrak Y\to\mathfrak Z$ deux morphismes de préschémas formels
  localement noethériens tels que $g$ soit de type fini et séparé~; alors, si
  $g\circ f$ est un morphisme fini, $f$ est un morphisme fini.
\end{enumerate}
\end{proposition}

\begin{proof}
(i) est trivial, et les autres assertions se ramènent aussitôt aux
propositions correspondantes pour les morphismes de préschémas usuels
(\textbf{II}, 6.1.5) à l'aide du critère \emph{a)} de
(\hyperref[III.4.8.1-fr]{4.8.1})~; nous laissons les détails au lecteur, sur le
modèle de (\textbf{I}, 10.13.5).
\end{proof}

\begin{corollary}[4.8.4]
\label{III.4.8.4-fr}
Sous les hypothèses de \textup{(\textbf{I}, 10.9.9)}, si $f$ est un morphisme fini, il en est de
même de son prolongement $\widehat f$ aux complétés.
\end{corollary}

\begin{corollary}[4.8.5]
\label{III.4.8.5-fr}
Si $\mathfrak X$ est un préschéma formel fini au-dessus de $\mathfrak Y$,
$f:\mathfrak X\to\mathfrak Y$ le morphisme structural, alors, pour tout ouvert
$U\subset\mathfrak Y$, $f^{-1}(U)$ est fini au-dessus de $U$.
\end{corollary}

\begin{proposition}[4.8.6]
\label{III.4.8.6-fr}
Si $f:\mathfrak X\to\mathfrak Y$ est un morphisme fini de préschémas formels
localement noethériens, $f_*(\mathcal O_{\mathfrak X})$ est une
$\mathcal O_{\mathfrak Y}$-Algèbre cohérente.
\end{proposition}

\begin{proof}
On peut considérer $f$ comme limite inductive d'un système inductif $(f_n)$ de
morphismes $f_{\boldsymbol{n}}:X_n\to Y_{\boldsymbol{n}}$~; montrons que les $f_n$ sont des morphismes finis
et $f_*(\mathcal O_{\mathfrak X})$ est isomorphe à la limite projective des
$(f_n)_*(\mathcal O_{X_n})$, ce qui établira notre assertion (\textbf{I}, 10.10.5). Il
suffit de se borner au cas où $\mathfrak Y=\operatorname{Spf}(B)$,
$\mathfrak X=\operatorname{Spf}(A)$, et de remarquer que si $\mathfrak R$ est
un idéal de définition de $B$ et $A$ un $B$-module de type fini,
$A/\mathfrak R^{n+1}A$ est un module de type fini sur
$B/\mathfrak R^{n+1}B$, et que $A$ est limite projective des
$A/\mathfrak R^{n+1}A$.
\end{proof}

Réciproquement~:

\begin{proposition}[4.8.7]
\label{III.4.8.7-fr}
Soient $\mathfrak Y$ un préschéma formel localement noethérien, $\mathcal A$
une $\mathcal O_{\mathfrak Y}$-Algèbre cohérente. Il existe un préschéma
formel $\mathfrak X$ fini au-dessus de $\mathfrak Y$, défini à un
$\mathfrak Y$-isomorphisme unique près, et tel que
$f_*(\mathcal O_{\mathfrak X})=\mathcal A$, $f:\mathfrak X\to\mathfrak Y$
étant le morphisme structural.
\end{proposition}

\begin{proof}
Soit $\mathcal K$ un Idéal de définition de $\mathfrak Y$, et posons
$Y_n=(\mathfrak Y,\mathcal O_{\mathfrak Y}/\mathcal K^{n+1})$ et
$\mathcal A_n=\mathcal A/\mathcal K^{n+1}\mathcal A$~;
il est clair que $\mathcal A_n$ est une $\mathcal O_{Y_n}$-Algèbre finie et
définit donc un

\end{proof}
\end{document}
